\section{Finite coding and bounded-formula Booleanization}
\label{sec:booleanization}

All later constructions depend on a modest but exact closure theorem: fixed finite data structures and
bounded logical operations can be represented by primitive-recursive functions, and their correctness can
be proved using quantifier-free induction.  This section states the construction in enough detail to be
implemented directly in the checked \texttt{PR} syntax.

\subsection{A concrete coding discipline}

The framework does not require a uniquely optimal G\"odel coding.  It requires one fixed coding with
primitive-recursive constructors, destructors, and size bounds.  The following choice is convenient.
For \(a,b\in\Nat\), let
\[
  \langle a,b\rangle = 2^a(2b+1).
\]
This code is positive and injective.  Given \(z>0\), the first component is the largest \(e\le z\) for
which \(2^e\) divides \(z\); the second is
\[
  \frac{z/2^e-1}{2}.
\]
Bounded division, remainder, exponentiation, divisibility, and bounded maximum are primitive recursive.
Thus pairing and both projections are primitive recursive on their intended range.  Values outside the
range may be sent to \(0\); every theorem is guarded by validity.

Encode the empty list by \(0\), and a cons cell by
\[
  \operatorname{cons}(a,t)=1+\langle a,t\rangle.
\]
The predicates \(\operatorname{isNil}\), \(\operatorname{head}\), and \(\operatorname{tail}\) are
primitive recursive.  Length and indexing are obtained by primitive recursion on a supplied step bound or,
for a valid list, by recursion on the code with a decreasing tail.  In object-theory developments it is
usually safer to carry an explicit length and prove that repeated tails reach \(0\) after that length.  This
avoids relying on a semantic recursion on the numeric size of an arbitrary malformed code.

Fixed-length vectors are more efficiently coded by mixed radix.  If coordinate \(i\) lies in
\([0,M_i]\), set
\[
  \operatorname{pack}(a_0,\ldots,a_{r-1})
   = (((a_0(M_1+1)+a_1)(M_2+1)+a_2)\cdots)(M_{r-1}+1)+a_{r-1}.
\]
For fixed \(r\), packing and each projection are primitive recursive.  The same representation later gives
a numerical rank for lexicographic descent.

\begin{lemma}[Finite-code closure]
For each fixed finite algebraic datatype generated from natural numbers by finite products, tagged sums,
and bounded lists, there are primitive-recursive validity, constructor, projection, equality, bounded fold,
and bounded-search operations.
\end{lemma}
\begin{proof}
Proceed by structural induction on the datatype description.  Products use pairing; tagged sums reserve a
finite tag field; a list of length at most \(B\) is iterated exactly \(B\) times, with a state containing the
current tail and accumulated result.  Every step is composition of previously constructed primitive-
recursive functions.  Since \(B\) is an input but the step function is fixed, the iteration is a primitive
recursion.  Validity is the conjunction of the tag, component, and termination checks.  Equality compares
all finitely visible fields and is a bounded fold.  No universal decoder is involved: a separate concrete
program is generated for each fixed datatype description.
\end{proof}

\subsection{Boolean primitives}

Write booleans as \(0\) and \(1\).  Primitive-recursive arithmetic supplies the following fixed functions:
\begin{align*}
  \operatorname{isZero}(x) &\in\Bool,\\
  \operatorname{eq}(x,y) &=
    \operatorname{isZero}\bigl((x\dotminus y)+(y\dotminus x)\bigr),\\
  \operatorname{lt}(x,y) &= \operatorname{isZero}(Sx\dotminus y),\\
  b\wedge_B c &= bc,\\
  b\vee_B c &= \min(1,b+c),\\
  \neg_B b &= 1\dotminus b,\\
  b\Rightarrow_B c &= (\neg_B b)\vee_B c.
\end{align*}
Here \(x\dotminus y\) is truncated subtraction.  Each operation is defined by primitive recursion and has
an object proof that its output is a boolean and that it reflects the corresponding decidable arithmetic
relation.

The important constructive point concerns implication.  In general intuitionistic logic, truth of
\(P\to Q\) is not captured by a two-valued truth table when \(P\) and \(Q\) are arbitrary propositions.
Here both propositions have explicit Boolean reflection.  From \(b_P=0\vee b_P=1\), a finite case split
shows
\[
  (b_P\Rightarrow_B b_Q)=1
  \quad\longleftrightarrow\quad
  (P\to Q).
\]
Thus Boolean implication is legitimate only after decidability has been constructed.

\subsection{Bounded formulas}

Fix an arity \(n\).  A bounded formula is generated by
\[
\begin{split}
 B ::= {}& t=u \mid t<u \mid B\wedge B \mid B\vee B \mid B\to B\\
        &\mid \forall z<t\;B(z) \mid \exists z<t\;B(z),
\end{split}
\]
where all terms are built from fixed primitive-recursive symbols.  A host-side compiler associates a
concrete characteristic term \(\chi_B(\vec x)\in\Bool\):
\begin{align*}
 \chi_{s=t} &= \operatorname{eq}(s,t), &
 \chi_{s<t} &= \operatorname{lt}(s,t),\\
 \chi_{C\wedge D} &= \chi_C\wedge_B\chi_D, &
 \chi_{C\vee D} &= \chi_C\vee_B\chi_D,\\
 \chi_{C\to D} &= \chi_C\Rightarrow_B\chi_D,\\
 \chi_{\forall z<t C} &= \prod_{z<t}\chi_C(z), &
 \chi_{\exists z<t C} &= \max_{z<t}\chi_C(z).
\end{align*}
The product and maximum are finite folds defined by primitive recursion.  Empty universal folds return
\(1\); empty existential folds return \(0\).

\begin{theorem}[Bounded Booleanization]
\label{thm:booleanization}
For every fixed bounded formula \(B(\vec x)\), the compiler produces a concrete primitive-recursive symbol
\(\chi_B\) and object-theory derivations of
\[
 \chi_B(\vec x)=1\;\leftrightarrow\;B(\vec x),
 \qquad
 \chi_B(\vec x)=0\;\leftrightarrow\;\neg B(\vec x).
\]
In particular, \(B\vee\neg B\) is derivable for this fixed bounded formula.
\end{theorem}

\begin{proof}
Induct on the syntax of \(B\).  Atomic equality and order use the verified Boolean arithmetic primitives.
Conjunction and disjunction follow by the truth-table equations and the induction hypotheses.  For
implication, first derive that each characteristic value is \(0\) or \(1\); perform the four finite cases;
then use the induction hypotheses to translate between Boolean values and formulas.

For a bounded universal quantifier, define
\[
  U(0,\vec x)=1,
  \qquad
  U(Sk,\vec x)=U(k,\vec x)\wedge_B\chi_C(k,\vec x).
\]
Use quantifier-free induction on \(k\) with invariant
\[
  U(k,\vec x)=1\;\leftrightarrow\;\forall z<k\;C(z,\vec x).
\]
The displayed bounded quantifier is an abbreviation expanded using the Boolean order test; after
Booleanization of its matrix, the invariant is equivalent to a quantifier-free equation or a fixed
first-order formula whose induction proof is assembled from the allowed quantifier-free invariant over the
state.  Operationally, it is often simplest to store both \(U(k)\) and a code of the bounded proof status in
the state, so the induction invariant itself is a conjunction of equations.  The successor step uses the
recursion equation for \(U\), the induction hypothesis, and the fact that
\(z<Sk\) splits decidably into \(z<k\) or \(z=k\).

The existential case is dual, using finite maximum or Boolean disjunction.  The zero-reflection statement
follows either from a parallel induction or from the fact that the characteristic is a boolean together
with the true-reflection statement.  Every recursion is attached to this particular formula; no formula
code is evaluated inside iPRA.
\end{proof}

\begin{remark}[Implementation form]
The cleanest Lean implementation separates two inductions.  A meta-level recursion on the finite syntax
\texttt{BForm} constructs a concrete \texttt{PR} term.  For each constructor it also constructs a
\(\Prf\,\PRAAx\) proof from previously generated proofs.  Bounded quantifiers invoke a generic fold lemma
proved once with \texttt{PRA.qfInd}.  This is proof-producing compilation, not reflection through an
untrusted semantic oracle.
\end{remark}

\subsection{Bounded minimization and complete enumeration}

Let \(P(x,i)\) be a bounded formula and \(B(x)\) a primitive-recursive bound.  Define a state with a
Boolean \(f\) and an index \(w\).  Initially \((f,w)=(0,0)\).  At stage \(i\):
\begin{itemize}
  \item if \(f=1\), leave the state unchanged;
  \item if \(f=0\) and \(\chi_P(x,i)=1\), set \((f,w)=(1,i)\);
  \item otherwise leave \(f=0\).
\end{itemize}
After \(B(x)\) steps, return \(w\) if \(f=1\), and a distinguished failure value otherwise.

\begin{theorem}[Least bounded search]
\label{thm:least-search}
There is a fixed primitive-recursive program \(\operatorname{least}_P(x,B)\) with an object proof of
\[
\begin{split}
 &\exists i<B\;P(x,i)
 \;\to\;
 P\bigl(x,\operatorname{least}_P(x,B)\bigr)
 \wedge \operatorname{least}_P(x,B)<B\\
 &\hspace{8em}\wedge
 \forall j<\operatorname{least}_P(x,B)\;\neg P(x,j),
\end{split}
\]
and an object proof that failure is equivalent to \(\forall i<B\;\neg P(x,i)\).
\end{theorem}

\begin{proof}
Use the quantifier-free invariant at stage \(t\):
\begin{enumerate}
  \item \(f\in\Bool\);
  \item if \(f=0\), then every \(j<t\) fails \(P\);
  \item if \(f=1\), then \(w<t\), \(P(x,w)\), and every \(j<w\) fails \(P\).
\end{enumerate}
Each bounded universal assertion is replaced by its Boolean fold, so the invariant is a conjunction of
term equations.  The successor proof is a decidable split on \(f\) and \(\chi_P(x,t)\).  At
\(t=B(x)\), the invariant gives both success and failure conclusions.  The leastness is therefore a proved
property of the concrete program, not an appeal to an external least-number principle.
\end{proof}

Complete bounded enumeration is similar.  Fold through \(i<B\), appending \(i\) exactly when
\(\chi_P(x,i)=1\).  The invariant says that the accumulated list is duplicate-free, every element is less
than \(t\) and satisfies \(P\), and every \(i<t\) satisfying \(P\) occurs in the list.  This yields a
canonical bridge from a height cutoff to a complete finite list, used in Section~\ref{sec:resources}.

\subsection{Why a universal evaluator is neither needed nor appropriate}

One could encode a bounded-formula syntax tree as a natural number and ask for a single arithmetic function
\(E(e,\vec x)\) evaluating every such code.  For a restricted language and an explicit depth bound, that
may be primitive recursive.  It is nevertheless the wrong default interface.  Application theorems use a
fixed formula, so host-side specialization produces a smaller concrete symbol, a simpler invariant, and a
proof that can be audited constructor by constructor.  More importantly, extending \(E\) to evaluate
arbitrary primitive-recursive program codes uniformly would conflict with diagonal arguments: there is no
single primitive-recursive function that is extensionally universal for all primitive-recursive unary
functions.  The staged design stays entirely on the safe side of this boundary.
